\chapter{Implementación}
\label{cap:implementacion}

En este capítulo se presentan los aspectos más relevantes de la implementación del simulador en C++17.

\section{Implementación del Cuerpo de Galois}
\label{sec:galois_field}

% (Se omiten los fundamentos matemáticos aquí; ver el capítulo de Fundamentos Matemáticos)
\subsection{Estructura de Datos}

\subsubsection{Atributos Estáticos}

La clase \texttt{GaloisField} utiliza atributos estáticos para compartir las tablas de precálculo entre todas las instancias del mismo cuerpo:

\begin{lstlisting}[style=CodigoCPP]
static vector<uint16_t> log_table;      // Tabla de logaritmos
static vector<uint16_t> exp_table;      // Tabla de exponentes
static int initialized_m;               // m del ultimo cuerpo inicializado
\end{lstlisting}

Esta estrategia reduce el consumo de memoria y las operaciones redundantes cuando se crean múltiples instancias del mismo cuerpo.

\subsubsection{Atributos de Instancia}

\begin{lstlisting}[style=CodigoCPP]
int m;                      // Grado del cuerpo
int size;                   // Numero de elementos (2^m)
int primitive_poly;         // Polinomio primitivo
\end{lstlisting}

\subsection{Inicialización y Precálculo}

El constructor inicializa el cuerpo y genera las tablas de logaritmos y exponentes:

\begin{lstlisting}[style=CodigoCPP]
GaloisField::GaloisField(int m, int primitive_poly) 
    : m(m), size(0), primitive_poly(primitive_poly) {
    if (m < 1 || m > 15) {
        throw std::invalid_argument("m debe estar entre 1 y 15");
    }
    if (primitive_poly <= 0) {
        throw std::invalid_argument(
            "El polinomio primitivo debe ser positivo");
    }
    
    size = 1 << m;  // size = 2^m

    // Solo se computa si es un nuevo cuerpo
    if (initialized_m != m) {
        log_table.assign(size, 0);
        exp_table.assign(2 * size, 0);

        int val = 1;
        for (int i = 0; i < size - 1; i++) {
            exp_table[i] = val;
            log_table[val] = i;
            val <<= 1;
            if (val & size) val ^= primitive_poly;
        }

        // Duplicar la tabla de exponentes para evitar modulo
        for (int i = size - 1; i < (size - 1) * 2; i++) {
            exp_table[i] = exp_table[i - (size - 1)];
        }
        initialized_m = m;
    }
}
\end{lstlisting}

\subsection{Operaciones Aritméticas}

Dado que las tablas exponencial (\texttt{exp\_table}) y logarítmica (\texttt{log\_table}) se calculan durante la inicialización, las operaciones aritméticas se resuelven en tiempo constante $O(1)$ mediante accesos a memoria, evitando cálculos algebraicos complejos durante la simulación.

\subsubsection{Adición} 
Se resuelve mediante una simple operación lógica XOR a nivel de bits (\texttt{a \textasciicircum{} b}), sin necesidad de acceder a las tablas.

\subsubsection{Multiplicación} 
Utiliza la propiedad de los logaritmos en cuerpos finitos:
\begin{equation}
a \cdot b = \alpha^{\log(a)} \cdot \alpha^{\log(b)} = \alpha^{\log(a) + \log(b)}
\end{equation}
A nivel de código, tras validar que ningún operando sea cero, esto se traduce en una única línea de C++:
\begin{lstlisting}[style=CodigoCPP]
return exp_table[log_table[a] + log_table[b]];
\end{lstlisting}
Al haber duplicado la tabla de exponentes en memoria durante la inicialización, se evita la costosa operación módulo $(2^m - 1)$ en cada ciclo de reloj.

\subsubsection{División} 
De forma análoga, la división se calcula restando los logaritmos:

\begin{equation}
\frac{a}{b} = \alpha^{\log(a) - \log(b)}
\end{equation}

En C++, se suma el tamaño máximo del cuerpo para garantizar que el índice sea siempre positivo sin necesidad de usar operaciones condicionales condicionales de salto:

\begin{lstlisting}[style=CodigoCPP]
return exp_table[log_table[a] - log_table[b] + (size - 1)];
\end{lstlisting}

\subsubsection{Inverso Multiplicativo y Potencia} 
Siguiendo la misma filosofía de optimización, el inverso de un elemento $a$ se obtiene directamente accediendo a \texttt{exp\_table[(size - 1) - log\_table[a]]}, y su potencia resolviendo \texttt{exp\_table[log\_table[a] * e \% (size - 1)]}.

\subsection{Gestión de Memoria y Optimizaciones}

El uso de tablas estáticas proporciona las siguientes ventajas:

\begin{itemize}
    \item \textbf{Precálculo único}: Si múltiples instancias utilizan el mismo $m$, las tablas se computan solo una vez
    \item \textbf{Acceso $O(1)$}: Todas las operaciones se reducen a búsquedas en tabla
    \item \textbf{Validación verificada}: El cuerpo \texttt{initialized\_m} garantiza consistencia
    \item \textbf{Complejidad espacial}: Para $\mathbb{GF}(2^m)$ aproximadamente $3 \times 2^{m+1}$ bytes. Para $m = 15$ (máximo soportado): $\approx 192$ KB
\end{itemize}

Para las operaciones aritméticas críticas (multiply, divide, power, inverse), la validación de los operandos se realiza mediante assert. Esta decisión de diseño se fundamenta en el rendimiento: al compilar el simulador en modo NDEBUG, el preprocesador elimina las aserciones, asegurando que no exista ninguna penalización (overhead) en tiempo de ejecución durante las simulaciones intensivas, mientras que mantiene la seguridad durante la fase de desarrollo y depuración.
\subsection{Ejemplo de Uso}

\begin{lstlisting}[style=CodigoCPP]
// Crear un cuerpo GF(2^4) con polinomio primitivo 0x13
GaloisField gf(4, 0x13);

// Operaciones
uint16_t a = 5;   // 0b0101 en binario
uint16_t b = 7;   // 0b0111 en binario

uint16_t suma = gf.add(a, b);           // 5 XOR 7 = 2 (0b0010)
uint16_t producto = gf.multiply(a, b);  // 5 * 7 = 8 (0b1000)
uint16_t cociente = gf.divide(a, b);    // 5 / 7 = 13 (0b1101)
uint16_t inverso = gf.inverse(b);       // 7^-1 = 6 (0b0110)
uint16_t potencia = gf.power(a, 3);     // 5^3 = 10 (0b1010)
\end{lstlisting}

\section{Implementación de Polinomios}
\label{sec:polynomial}

\subsection{Estructura de datos}

La clase \texttt{Polynomial} representa polinomios con coeficientes en $\mathbb{GF}(2^m)$ usando:

\begin{lstlisting}[style=CodigoCPP]
// coef[0] + coef[1]x + coef[2]x^2 + ...
vector<uint16_t> coef;
GaloisField gf; // cuerpo asociado
\end{lstlisting}

\texttt{trim()} mantiene la representación eliminando ceros finales.

\subsection{Constructor y normalización}

\begin{lstlisting}[style=CodigoCPP]
Polynomial::Polynomial(const GaloisField &galoisField,const vector<uint16_t> &coefficients)
    : coef(coefficients), gf(galoisField) {
    trim();
}
\end{lstlisting}

\subsection{Operaciones implementadas}

\subsubsection{Acceso y modificación}

\begin{lstlisting}[style=CodigoCPP]
int Polynomial::getDegree() const { /* devuelve coef.size()-1 o -1 si vacio */ }
uint16_t Polynomial::getCoef(int degree) const { /* devuelve 0 si fuera de rango */ }
int Polynomial::setCoef(int degree, uint16_t value) { /* redimensiona y trim */ }
\end{lstlisting}

\subsection{Operaciones Algebraicas}

La clase implementa las operaciones algebraicas fundamentales operando internamente sobre los vectores de coeficientes en $\mathbb{GF}(2^m)$.

\begin{itemize}
    \item \textbf{Suma:} Se realiza coeficiente a coeficiente alineando los grados de ambos polinomios. Dado que la aritmética subyacente es $GF(2^m)$, esta operación tiene una complejidad $O(\max(\deg(p_1), \deg(p_2)))$ y equivale a un XOR vectorial.
    \item \textbf{Multiplicación:} Se ha implementado el método estándar. Para cada par de coeficientes, se utiliza la multiplicación del cuerpo de Galois, resultando en un polinomio de grado $\deg(p_1) + \deg(p_2)$ con una complejidad computacional de $O(\deg(p_1) \cdot \deg(p_2))$.
    \item \textbf{Evaluación (Método de Horner):} Para evaluar un polinomio en un elemento $x$ del cuerpo (operación crítica durante el cálculo de síndromes y la búsqueda de Chien), se ha implementado el método de Horner. Esta técnica evita el costoso cálculo explícito de potencias, reduciendo la evaluación a una serie anidada de multiplicaciones y sumas:
    \begin{equation}
    P(x) = (\dots((a_n x + a_{n-1})x + a_{n-2})x + \dots )x + a_0
    \end{equation}
    Esto garantiza que cualquier polinomio se evalúe en tiempo lineal $O(n)$ realizando exactamente $n$ multiplicaciones y $n$ sumas en el cuerpo finito.
\end{itemize}

\subsubsection{Trim y utilidades}

\begin{lstlisting}[style=CodigoCPP]
void Polynomial::trim() {
    while (!coef.empty() && coef.back() == 0) coef.pop_back();
}
\end{lstlisting}

\subsection{Decisiones de diseño}

\begin{itemize}
  \item Algoritmo de multiplicación simple: suficiente para los grados usados en este proyecto; alternativas (Karatsuba) solo si se requiere rendimiento en polinomios grandes.
  \item Uso de \texttt{uint16\_t} para coeficientes acorde a \texttt{GaloisField} actual ($m\leq15$).
\end{itemize}

\section{Implementación de PackedBit}
\label{sec:packedbit}

\subsection{Estructura de datos}

La clase \texttt{PackedBit} empaqueta una secuencia de bits lógicos dentro de un vector de bloques de 64 bits (\texttt{vector<uint64\_t>}). Esta representación reduce el coste de memoria frente a almacenar cada bit como un entero independiente y permite aplicar XOR y desplazamientos a nivel de palabra completa.

En el contexto del codificador BCH, \texttt{PackedBit} se usa para representar tanto el mensaje extendido como el polinomio generador en forma binaria, acelerando la división polinómica por XOR desplazado.


\subsection{Acceso y modificación de bits}

\begin{lstlisting}[style=CodigoCPP]
uint16_t get(size_t index) const;
void set(size_t index, uint16_t value);
\end{lstlisting}

\texttt{get()} devuelve 0 cuando el índice está fuera de rango, y \texttt{set()} ignora escrituras fuera del tamaño lógico del objeto.
\subsection{Empaquetado y Desempaquetado de Bits}

La transición entre un elemento \texttt{std::vector\textless uint16\_t\textgreater} a \texttt{uint64\_t} se gestiona de forma transparente en el constructor y en el método de exportación \texttt{toVector()}. El constructor calcula el número de bloques de 64 bits necesarios (\texttt{blockCount} = (\texttt{bitCount} + 63) / 64) y aplica máscaras de bits (\texttt{1ULL} $\ll$ \texttt{offset}) para condensar la información. A la inversa, \texttt{toVector()} extrae los bits lógicos procesando bloques completos y gestionando el residuo (\texttt{bitCount} \% 64) para devolver la palabra código al entorno de simulación.

\subsection{Desplazamiento y XOR combinado}

La operación \texttt{xorShifted()} es la más importante para el codificador BCH, porque permite simular la resta polinómica mediante XOR del generador desplazado.

\begin{lstlisting}[style=CodigoCPP]
void xorShifted(const PackedBit &gx, size_t shift_bits) {
    if (shift_bits >= bitCount || gx.data.empty()) return;

    size_t block_shift = shift_bits / 64;
    size_t bit_shift = shift_bits % 64;
    uint64_t carry = 0, next_carry = 0;
    
    size_t dst;
    for (size_t i = 0; i < gx.data.size(); ++i) {
        dst = i + block_shift;
        if (dst >= data.size()) break;

        uint64_t word = gx.data[i];
        next_carry = (bit_shift == 0) ? 0 : (word >> (64 - bit_shift));
        data[dst] ^= (word << bit_shift) | carry;
        carry = next_carry;
    }

    dst = gx.data.size() + block_shift;
    if (carry && dst < data.size()) {
        data[dst] ^= carry;
    }

    if (bitCount % 64 != 0 && !data.empty()) {
        uint64_t keep_mask = (1ULL << (bitCount % 64)) - 1ULL;
        data.back() &= keep_mask;
    }
}
\end{lstlisting}

\subsection{Decisiones de diseño}

\begin{itemize}
    \item Se usa \texttt{uint64\_t} como unidad base para aprovechar operaciones de la CPU.
    \item \texttt{bitCount} conserva el tamaño lógico original del conjunto de bits.
    \item Se prioriza la velocidad de operaciones binarias frente a una abstracción más general.
    \item La versión optimizada de \texttt{toVector()} reduce el coste de conversión cuando el codificador necesita recuperar los bits residuales.
\end{itemize}

\subsection{Relación con BCH\_Codec}

\texttt{PackedBit} no es solo una estructura auxiliar, sino la representación binaria que permite que \texttt{BCH\_Codec} realice la división del mensaje por el polinomio generador mediante desplazamientos y XOR sobre bloques de 64 bits. En otras palabras, convierte una operación algebraica sobre polinomios en una secuencia de operaciones bit a bit mucho más eficiente.

\section{Implementación de \texttt{BCH\_Codec}}
\label{sec:bch_codec}

En esta sección se presenta la implementación completa de la clase \texttt{BCH\_Codec}, que concentra tanto la generación del polinomio generador como la codificación, la detección de errores y la corrección de bloques recibidos. A diferencia de las clases anteriores, aquí no se implementa una sola operación aislada, sino el flujo completo del codificador/decodificador BCH.

\subsection{Estructura de datos}

Los atributos relevantes de la clase \texttt{BCH\_Codec} son:

\begin{lstlisting}[style=CodigoCPP]
int m;                // Grado del cuerpo
int t;                // Capacidad de correccion
int n;                // Longitud del codigo (2^m - 1)
int k;                // Longitud del mensaje (n - grado(generator))
GaloisField gf;       // cuerpo asociado
Polynomial generator; // Polinomio generador
\end{lstlisting}

Además, la clase mantiene un método privado \texttt{computeGeneratorPolynomial()} encargado de construir el polinomio generador a partir de los polinomios mínimos necesarios para el BCH elegido.

\subsection{Constructor y cálculo del generador}

El constructor comprueba que $1\le m\le 15$ y $1 \le 2t < n$. Una vez validados los parámetros, inicializa el cuerpo de Galois, calcula \texttt{n = $2^m - 1$} y determina \texttt{k} a partir del grado del polinomio generador. La construcción del generador se delega a \texttt{computeGeneratorPolynomial()}, que recorre las raíces conjugadas y acumula los polinomios mínimos correspondientes.

\begin{lstlisting}[style=CodigoCPP]
BCH_Codec::BCH_Codec(int m, int t, int primitive_poly)
    : m(m), t(t), gf(m, primitive_poly), generator(gf, {}) {
    if (m < 1 || m > 15) {
        throw std::invalid_argument("Invalid BCH parameters: m must be between 1 and 15");
    }

    n = (1 << m) - 1;

    if (t <= 0 || 2LL * t >= n) {
        throw std::invalid_argument("Invalid BCH parameters: t must satisfy 1 <= 2t < n");
    }

    computeGeneratorPolynomial();
    k = n - generator.getDegree();
}
\end{lstlisting}

\begin{lstlisting}[style=CodigoCPP]
void BCH_Codec::computeGeneratorPolynomial() {
    generator = Polynomial(gf, std::vector<uint16_t>{1});
    std::unordered_set<uint16_t> processed_exp;

    for (uint16_t i = 1; i <= 2 * t; i++) {
        if (processed_exp.count(i)) continue;

        Polynomial mi(gf, {1});
        uint16_t currExp = i;

        do {
            uint16_t root = gf.power(2, currExp);
            mi = mi * Polynomial(gf, {root, 1});
            processed_exp.insert(currExp);
            currExp = (currExp * 2) % n;
        } while (currExp != i);

        generator = generator * mi;
    }
}
\end{lstlisting}

\subsection{Codificación sistemática (versión empaquetada)}

La versión optimizada utilizada en el simulador convierte el mensaje y el generador en \texttt{PackedBit} y realiza la división polinómica mediante XOR desplazado. Esta codificación es sistemática: los bits de información se conservan y los bits de paridad se añaden al principio de la palabra de código.

Las etapas principales son:

\begin{enumerate}
  \item Extender el mensaje añadiendo $g$ ceros a la izquierda, donde $g=\deg(generator)$.
  \item Construir \texttt{PackedBit} para el vector extendido y para el polinomio generador.
  \item Iterar desde el bit más significativo; si el bit es 1, aplicar \texttt{xorShifted} con el generador alineado.
  \item Convertir el buffer resultante a vector y extraer los primeros $g$ bits como paridad. Concatenar paridad y mensaje.
\end{enumerate}

\begin{lstlisting}[style=CodigoCPP]
std::vector<uint16_t> BCH_Codec::encode(const std::vector<uint16_t>& message) {
    if (message.size() != static_cast<size_t>(k))
        throw std::invalid_argument("Message length must be exactly k");

    int g = generator.getDegree();

    // 1. Mensaje extendido: m(x) * x^g
    std::vector<uint16_t> buffer(g, 0);
    buffer.insert(buffer.end(), message.begin(), message.end());
    PackedBit binaryBuffer(buffer);

    // 2. Generador en forma binaria
    std::vector<uint16_t> gen_bits(g+1, 0);
    for (int deg = 0; deg <= g; ++deg) gen_bits[deg] = generator.getCoef(deg);
    PackedBit gen_packed(gen_bits);

    // 3. Division polinomica: si MSB = 1, XOR del generador alineado
    for (int i = buffer.size() - 1; i >= g; --i) {
        if (binaryBuffer.get(i)) {
            size_t shift = (i - g);
            binaryBuffer.xorShifted(gen_packed, shift);
        }
    }

    // 4. Extraer resto y formar palabra sistematica
    buffer = binaryBuffer.toVector();
    std::vector<uint16_t> parity(buffer.begin(), buffer.begin() + g);
    parity.insert(parity.end(), message.begin(), message.end());
    return parity;
}
\end{lstlisting}

\subsection{Detección de errores mediante síndromes}

Para el proceso de decodificación, el primer paso consiste en calcular el vector de síndromes. Si todos son cero, el bloque recibido es válido; en caso contrario, se detecta la presencia de errores.

\begin{lstlisting}[style=CodigoCPP]
std::vector<uint16_t> BCH_Codec::syndrome(const std::vector<uint16_t> &received, bool &error) {
    Polynomial r(gf, received);
    std::vector<uint16_t> synd(2 * t + 1, 0);
    error = false;
    for(int i=1; i<=2*t; i++){
        if (i%2 == 0){
            synd[i] = gf.power(synd[i/2],2);
        } else{
            uint16_t root = gf.power(2, i);
            synd[i] = r.evaluate(root);
        }
        if (synd[i] != 0 and !error){
            error = true;
        }
    }
    return synd;
}
\end{lstlisting}

La implementación aprovecha la propiedad de los síndromes en cuerpos binarios: los índices pares se obtienen por cuadrado del síndrome anterior, lo que reduce el número de evaluaciones directas del polinomio recibido.

\subsection{Localización de errores con Berlekamp-Massey}

Una vez detectada la presencia de errores, la clase estima el polinomio localizador de errores mediante el algoritmo de Berlekamp-Massey. 

En lugar de presentar el código fuente íntegro, a continuación se detalla la lógica algorítmica implementada:

\begin{enumerate}
    \item \textbf{Inicialización:} Se definen los polinomios $\sigma(x) = 1$ y un polinomio auxiliar $B(x) = 1$. Se establece la longitud estimada de errores $L = 0$.
    \item \textbf{Iteración principal:} Para cada síndrome $S_i$ (desde $i = 1$ hasta $2t$):
    \begin{itemize}
        \item Se calcula la discrepancia $d$, evaluando si el síndrome actual es predecible mediante el polinomio $\sigma(x)$ actual.
        \item Si $d = 0$, el polinomio actual es válido para este paso; se incrementa el desplazamiento del polinomio auxiliar $B(x)$.
        \item Si $d \neq 0$, se actualiza $\sigma(x)$ restándole el polinomio auxiliar ajustado por $d$. Si se supera la cota de longitud ($2L \leq i - 1$), se actualiza la estimación de errores $L$ y se guarda el polinomio previo como nuevo auxiliar $B(x)$.
    \end{itemize}
\end{enumerate}

\textbf{Decisión de diseño en C++:} A nivel de implementación, programar este algoritmo instanciando objetos de la clase \texttt{Polynomial} en cada iteración generaría un cuello de botella inaceptable. Por ello, el algoritmo se ha reescrito para operar directamente sobre \texttt{std::vector<uint16\_t>}. Además, se ha asignado el tamaño máximo teórico de los polinomios ($2t+2$) de forma estática al inicio de la función, eliminando por completo las costosas reasignaciones de memoria (\texttt{.resize()}) durante el bucle.

\subsection{Decodificación}

La decodificación combina la detección de errores, la estimación del polinomio localizador y la búsqueda de raíces de ese polinomio mediante la búsqueda de Chien. El flujo de ejecución implementado es el siguiente:

\begin{enumerate}
    \item \textbf{Cálculo de Síndromes:} Se procesa el bloque recibido. Si todos los síndromes son nulos, se asume que no hay errores, finalizando la ejecución prematuramente (camino rápido).
    \item \textbf{Localización (Berlekamp-Massey):} Si hay errores, se genera el polinomio localizador $\sigma(x)$. Si el grado de este polinomio supera la capacidad de corrección ($L > t$), se marca el bloque como incorregible (CWER) y se aborta para ahorrar tiempo de cómputo.
    \item \textbf{Búsqueda de Raíces (Chien Search):} Se evalúa $\sigma(x)$ iterativamente para todos los elementos del cuerpo. 
    \item \textbf{Corrección:} Cada vez que la evaluación resulta en cero (indicando una raíz), se localiza la posición física del error en el bloque original y se invierte dicho bit mediante una operación lógica XOR (\texttt{corrected[i] \textasciicircum{}= 1}).
    \item \textbf{Comprobación de Integridad:} Finalmente, se verifica que el número de raíces encontradas en el canal coincida exactamente con el grado del polinomio localizador $L$. De no ser así, indica que el canal superó los límites teóricos del código, devolviendo un fallo de corrección.
\end{enumerate}

Este enfoque modular garantiza que el sistema solo invierta ciclos de reloj en la costosa búsqueda de Chien cuando es matemáticamente seguro que los errores son corregibles.

\subsection{Comentarios de diseño}

\begin{itemize}
  \item \textbf{Rendimiento}: \texttt{PackedBit} reduce el número de operaciones XOR.
  \item \textbf{Reutilización}: El cálculo del generador y la decodificación están encapsulados en métodos privados y públicos separados, lo que hace más legible el flujo de la clase.
\end{itemize}

\section{Implementación del Canal de Ruido}
\label{sec:channel}

El componente \texttt{Channel} modela los principales canales de ruido usados en las simulaciones: el canal binario simétrico (BSC) y el canal AWGN con decisión dura. La clase encapsula un generador de números aleatorios \texttt{std::mt19937} y expone métodos claros para aplicar cada modelo de ruido.

\subsection{Interfaz pública}

La interfaz de \texttt{Channel} se resume en:

\begin{lstlisting}[style=CodigoCPP]
class Channel{
private:
    std::mt19937 rng; // Generador de numeros aleatorios
public:
    Channel();
    std::vector<uint16_t> applyBSC(const std::vector<uint16_t>& bits, double bit_error_rate);
    std::vector<double> transmitAWGN(const std::vector<uint16_t>& bits, double es_n0_db);
    std::vector<uint16_t> applyAWGNHardDecision(const std::vector<uint16_t>& bits, double es_n0_db);
};
\end{lstlisting}

\subsection{Constructor}

El constructor inicializa \texttt{rng} con una semilla no determinista tomada de \texttt{std::random\_device} para obtener trayectorias de ruido reproducibles solo si se sustituye la semilla.

\begin{lstlisting}[style=CodigoCPP]
Channel::Channel() {
    std::random_device rd;
    rng = std::mt19937(rd());
}
\end{lstlisting}

\subsection{Canal Binario Simétrico (BSC)}

\texttt{applyBSC} aplica fallos independientes por bit con probabilidad \texttt{bit\_error\_rate} usando una distribución de Bernoulli. Valida el parámetro de probabilidad y lanza \texttt{std::invalid\_argument} si está fuera de rango.

\begin{lstlisting}[style=CodigoCPP]
std::vector<uint16_t> Channel::applyBSC(const std::vector<uint16_t>& bits, double bit_error_rate) {
    if (bit_error_rate < 0.0 || bit_error_rate > 1.0) {
        throw std::invalid_argument("bit_error_rate must be between 0 and 1");
    }

    std::bernoulli_distribution flip(bit_error_rate);
    std::vector<uint16_t> noisy_bits = bits;
    for (uint16_t& bit : noisy_bits) {
        if (flip(rng)) bit ^= 1;
    }

    return noisy_bits;
}
\end{lstlisting}

\subsection{Canal AWGN}

Para señalización binaria antipodal se mapea cada bit en símbolo $+1$ o $-1$ y se añade ruido Gaussiano con desviación estándar calculada a partir de \texttt{$E_s/N_0$}. La función auxiliar \texttt{dbToLinear} convierte dB a escala lineal.

\begin{lstlisting}[style=CodigoCPP]
double dbToLinear(double db_value) {
    return std::pow(10.0, db_value / 10.0);
}

std::vector<double> Channel::transmitAWGN(const std::vector<uint16_t>& bits, double es_n0_db) {
    const double es_n0_linear = dbToLinear(es_n0_db);
    const double sigma = std::sqrt(1.0 / (2.0 * es_n0_linear));

    std::normal_distribution<double> noise(0.0, sigma);
    std::vector<double> samples; samples.reserve(bits.size());

    for (uint16_t bit : bits) {
        const double symbol = (bit == 0) ? 1.0 : -1.0;
        samples.push_back(symbol + noise(rng));
    }

    return samples;
}
\end{lstlisting}

\subsection{Decisión dura sobre AWGN}

La función \texttt{applyAWGNHardDecision} realiza la transmisión AWGN y aplica una decisión dura: si la muestra es $\ge 0$ se interpreta como 0, en caso contrario como 1.

\begin{lstlisting}[style=CodigoCPP]
std::vector<uint16_t> Channel::applyAWGNHardDecision(const std::vector<uint16_t>& bits, double es_n0_db) {
    std::vector<double> samples = transmitAWGN(bits, es_n0_db);
    std::vector<uint16_t> hard_bits; hard_bits.reserve(samples.size());
    for (double sample : samples) hard_bits.push_back(sample >= 0.0 ? 0u : 1u);
    return hard_bits;
}
\end{lstlisting}

\subsection{Decisiones de diseño y uso}

\begin{itemize}
  \item \textbf{Separación de responsabilidades}: el canal ofrece modelos de canal (BSC, AWGN) como funciones independientes, facilitando su comparación experimental.
  \item \textbf{Reproducibilidad}: la semilla de \texttt{rng} se extrae de \texttt{random\_device} por defecto; para tests reproducibles puede sustituirse por una semilla fija.
\end{itemize}

\section{Arquitectura del Simulador y Motor de Ejecución}
\label{sec:bch_simulator}

La clase \texttt{BCH\_Simulator} actúa como el motor principal del proyecto. Su responsabilidad es orquestar las instancias del codificador/decodificador (\texttt{BCH\_Codec}) y del canal ruidoso (\texttt{Channel}), gestionar el barrido de relaciones señal a ruido ($E_b/N_0$), y recolectar las métricas de rendimiento para su posterior análisis.

\subsection{Estructuras de Configuración y Resultados}

Para garantizar la modularidad, los parámetros de entrada se encapsulan en la estructura \texttt{SimulationConfig}, la cual define los parámetros del código ($m, t$), el rango de simulación, y los criterios de parada estadística. De forma análoga, la estructura \texttt{PointResults} consolida las métricas de salida por cada punto de $E_b/N_0$.

\begin{lstlisting}[style=CodigoCPP]
struct PointResults { 
    double ebno_db; 
    double ber;             // Bit Error Rate (Codificado)
    double ber_uncoded;     // Bit Error Rate (Sin codificar)
    double cwer;            // Codeword Error Rate
    long codewords_simulated; 
    double avg_enc_us;      // Latencia media de codificacion
    double avg_dec_us;      // Latencia media de decodificacion
};

struct SimulationConfig { 
    int m = 7; 
    int t = 10; 
    double ebno_min = 0.0; 
    double ebno_max = 12.0; 
    double step = 1.0; 
    int min_codeword_errors = 100; // Criterio de parada Monte Carlo
    long max_codewords = 1000000; 
    bool use_file = false;
    std::string input_path = ""; 
};
\end{lstlisting}

\subsection{Motor de Simulación y Paralelización (OpenMP)}

El núcleo computacional del simulador reside en el método \texttt{simulatePoint}, encargado de generar las transmisiones para un valor específico de $E_b/N_0$. Dado que las simulaciones de Monte Carlo para tasas de error bajas (BER $< 10^{-5}$) requieren procesar millones de bloques, este método ha sido paralelizado utilizando OpenMP.

\subsubsection{Gestión de Concurrencia y Condición de Parada}

Para implementar este motor de simulación, la paralelización no se aborda mediante hilos del sistema operativo tradicionales (\textit{pthreads} o \texttt{std::thread}), sino utilizando directivas de compilador \textbf{OpenMP}. Esto permite un modelo \textit{Fork-Join} de muy baja sobrecarga.

El diseño concurrente del bucle principal se centra en torno a tres decisiones clave:

\begin{enumerate}
    \item \textbf{Aislamiento del Estado (\textit{Thread-Local Storage}):} 
    Para evitar condiciones de carrera, cada hilo del procesador instancia su propia copia local del códec BCH y del modelo de canal. Fundamentalmente, cada hilo recibe un motor de generación de números pseudoaleatorios (\texttt{std::mt19937}) inicializado con una semilla independiente, garantizando la incorrelación del ruido.

    \item \textbf{Balanceo de Carga Dinámico (\texttt{schedule(dynamic)}):} 
    El tiempo de decodificación no es constante; una palabra código con muchos errores requerirá ejecutar la costosa búsqueda de Chien, mientras que una palabra limpia se procesa de inmediato. Al emplear planificación dinámica, OpenMP asigna nuevas iteraciones a los hilos conforme terminan las anteriores, evitando que los hilos rápidos queden bloqueados esperando a los lentos.

    \item \textbf{Reducciones vs. Operaciones Atómicas:} 
    La actualización del marcador global de bits erróneos es una operación de altísima frecuencia. Si todos los hilos intentaran modificar una variable global simultáneamente usando exclusión mutua o \texttt{\#pragma omp atomic}, se produciría un severo cuello de botella en el hardware conocido como contención de caché (\textit{cache line bouncing}). Para evitar esto, el simulador hace uso de la cláusula \texttt{reduction}: OpenMP dota a cada hilo de una copia privada temporal para acumular sus errores y, al finalizar el bucle, el propio \textit{framework} obtiene todos los subtotales de forma segura. 
    Las lecturas y escrituras atómicas (\texttt{atomic read / write}) se han reservado únicamente para la variable compartida de tramas erróneas, permitiendo que todos los hilos monitoricen la condición de parada sin impactar el rendimiento global.
\end{enumerate}

\subsection{Decisiones de Diseño Adicionales}

\begin{itemize}
    \item \textbf{Inyección de Tráfico Real}: Además del generador pseudoaleatorio, el simulador permite inyectar archivos binarios reales (\textit{files}) como fuente de datos. Esto no solo valida la total transparencia del sistema ante cualquier secuencia de bits, sino que permite extraer archivos corruptos post-simulación para analizar cualitativamente el impacto de los errores residuales (por ejemplo, evaluando la degradación visual en archivos de imagen). 
    \item \textbf{Gestión Eficiente de Memoria}: Para evitar cuellos de botella en la memoria RAM al cargar estos archivos de gran tamaño, se emplea la librería estándar \texttt{<filesystem>} (introducida en C++17) para pre-calcular el tamaño exacto del archivo en bytes. Esto permite realizar una reserva de memoria estática (\texttt{reserve}) de un solo golpe antes de iniciar la lectura binaria. Además, la escritura de los bits corregidos por parte de los hilos de OpenMP se realiza de forma quirúrgica calculando índices de memoria directos (\texttt{base\_index}), eludiendo el uso de secciones críticas que paralizarían la ejecución paralela.
    \item \textbf{Métricas de Latencia}: La recolección de los tiempos de codificación y decodificación (\texttt{avg\_enc\_us}, \texttt{avg\_dec\_us}) se realiza mediante la librería \texttt{<chrono>}.
    \item \textbf{Exportación}: Todas las métricas se vuelcan automáticamente en formato \texttt{CSV}, permitiendo su uso directo en herramientas de postprocesado como Python (Matplotlib/Pandas) para la generación de las curvas de rendimiento abordadas en el siguiente capítulo.
\end{itemize}